\documentclass[11pt]{article}
\usepackage[utf8]{inputenc}
\usepackage[margin=1in]{geometry}
\usepackage{amsmath}
\usepackage{graphicx}
\usepackage{booktabs}
\usepackage{hyperref}
\usepackage{natbib}
\usepackage{setspace}
\onehalfspacing

\title{The Breath Shape Controls Intonation of Mouse Ultrasonic Vocalizations}
\author{Author A$^{1}$, Author B$^{1,2}$, Author C$^{2}$ \\
\small $^{1}$Department of Neuroscience, University of Research \\
\small $^{2}$Department of Otolaryngology, Medical Center}
\date{}

\begin{document}
\maketitle

\begin{abstract}
Mouse ultrasonic vocalizations (USVs) comprise a rich repertoire of syllable types defined by their pitch contours, but the mechanisms generating this diversity are poorly understood. The prevailing view holds that dynamic changes in laryngeal tension are the sole mechanism for pitch modulation. Here we demonstrate that modulation of the expiratory breath waveform is a primary mechanism controlling pitch patterns (intonation) in 6 of 10 USV syllable types. Using simultaneous whole-body plethysmography and microphone recordings in male mice exposed to female urine, we found that instantaneous pitch was positively correlated with expiratory airflow changes in the majority of syllable types (positive intonation), while a minority showed anticorrelated pitch-airflow relationships (negative intonation, presumably laryngeal). Electromyography recordings revealed that re-activation of the diaphragm (an inspiratory muscle) during expiration creates the breath waveform modulations underlying positive intonation. Optogenetic activation of the intermediate Reticular Oscillator (iRO), a brainstem vocalization central pattern generator, was sufficient to produce most endogenous USV syllable types with primarily positive intonation patterns. These results establish a two-mechanism model of USV intonation: a breath-driven mechanism mediated by the iRO that controls the majority of syllable types, and a laryngeal mechanism that generates the remaining types. The two mechanisms can combine within single vocalizations to produce complex pitch patterns.
\end{abstract}

\section{Introduction}

Mouse ultrasonic vocalizations (USVs) are a primary model system for studying the neural control of vocal communication, with applications to autism spectrum disorder, Fragile X syndrome, and other neurodevelopmental conditions that affect social communication \citep{holy2005,portfors2007,scattoni2009}. Male mice produce complex sequences of USVs in social contexts, particularly during courtship interactions with females. These vocalizations span 40--120 kHz and are organized into approximately 10 distinct syllable types defined by their pitch contours (frequency modulation patterns): flat, down FM, up FM, chevron, step up, step down, short, complex, multi-step, and two-syllable calls \citep{grimsley2011}.

The mechanisms by which mice generate this diversity of pitch patterns are central to understanding how USVs encode social information. The prevailing view in mammalian vocal physiology is that pitch modulation is controlled primarily by changes in laryngeal muscle tension, particularly the cricothyroid muscle, which stretches the vocal folds to increase pitch \citep{titze1994,riede2011}. Expiratory airflow is generally thought to modulate loudness (amplitude) but not pitch.

However, this view faces a paradox. Experiments in which air is injected below the larynx of anesthetized rodents demonstrate that increased subglottal pressure increases pitch \citep{riede2013}. Yet in natural vocalizations, the correlation between breath airflow and pitch is often weak or absent \citep{sirotin2014}. This discrepancy suggests that the relationship between airflow and pitch in natural USVs is more complex than simple pressure-pitch coupling.

The intermediate Reticular Oscillator (iRO) is a recently described brainstem circuit that functions as a central pattern generator (CPG) for vocalization in neonatal mice \citep{wei2022}. The iRO coordinates breathing rhythms and phonation through glutamatergic (Vglut2+) neurons expressing proenkephalin (Penk+) located in the ventral intermediate reticular tract, medial to the compact nucleus ambiguus. Whether the iRO plays a role in adult USV production and, specifically, in generating the breath waveform modulations that might control intonation has not been investigated.

Here we combine whole-body plethysmography, microphone recording, electromyography of respiratory and laryngeal muscles, and optogenetic circuit manipulation to demonstrate that breath shape---specifically, re-activation of the diaphragm during expiration---is a primary mechanism controlling USV intonation through the iRO brainstem circuit.

\section{Methods}

\subsection{Animals and Behavioral Setup}

Male C57BL/6 mice ($n = 6$) were placed in modified whole-body plethysmography chambers equipped with ultrasonic microphones for simultaneous recording of breathing and USVs. Fresh female urine was introduced to elicit courtship vocalizations. Recordings were performed during the first 5--10 minutes of urine exposure, when vocalization rates peak.

\subsection{Vocalization Classification}

USVs were defined as narrow-band ultrasonic sounds (40--120 kHz) occurring during single breaths. Syllable types were classified according to standard taxonomy based on pitch contour shape \citep{grimsley2011}: flat, down FM, up FM, chevron, step up, step down, short, complex, multi-step, and two-syllable.

\subsection{Intonation Analysis}

For each vocalization, the instantaneous correlation between expiratory airflow (from plethysmography) and pitch (from spectrogram) was computed. Positive intonation was defined as $r > 0$ (pitch tracks airflow), negative intonation as $r < 0$ (pitch anticorrelates with airflow). The magnitude and sign of this correlation for each syllable type determined the predominant intonation mechanism.

\subsection{Electromyography}

Chronic EMG electrodes were implanted in the diaphragm (inspiratory muscle), thyroarytenoid (laryngeal adductor), and cricothyroid (laryngeal tensor) muscles. EMG was recorded simultaneously with plethysmography and USV microphone recordings to identify muscle activation patterns during vocalization.

\subsection{Optogenetic Experiments}

iRO neurons were targeted for optogenetic activation using intersectional genetics. Penk-Cre;Vglut2-Flp mice received bilateral injections of AAV CreON-FlpON-ChR2::YFP into the iRO region. Optical fibers were implanted over the injection sites. Light stimulation at various frequencies (e.g., 10 Hz) was delivered during plethysmography and USV recording to assess the sufficiency of iRO activation for producing vocalizations.

\subsection{Statistical Analysis}

Breath parameter comparisons between USV and non-USV breaths used paired $t$-tests. Intonation correlations were assessed using Pearson correlation and one-way ANOVA with Sidak's post-hoc correction across syllable types.

\section{Results}

\subsection{Basic Vocalization and Breathing Parameters}

USVs occurred predominantly during expiration, with 88\% of breaths containing a single syllable (Table~\ref{tab:basic}). The mean instantaneous frequency of USV breaths was 7.5 Hz, slightly slower than neighboring non-USV sniff breaths (8.5--10 Hz; paired $t$-test $p = 0.03$). Inspiratory and expiratory timing parameters were similar between USV and non-USV breaths (Ti: $p = 0.40$; Te: $p = 0.18$; Ti/Te ratio: $p = 0.25$), but peak inspiratory flow was significantly larger for USV breaths ($p = 0.01$), suggesting that mice take deeper breaths before vocalizing.

\begin{table}[h]
\centering
\caption{Basic vocalization and breathing parameters ($n = 6$ mice).}
\label{tab:basic}
\begin{tabular}{ll}
\toprule
Parameter & Value \\
\midrule
Vocalization frequency range & 40--120 kHz \\
Typical fundamental frequency & $\sim$75 kHz \\
Peak vocalization rate & $\sim$4 events/second \\
Mean USV breath frequency & 7.5 Hz \\
Sniff episode frequency & 8.5--10 Hz \\
Syllables per breath & 88\% single syllable \\
USV onset & Biased to early expiration \\
\bottomrule
\end{tabular}
\end{table}

\subsection{Syllable Type Distribution}

The 10 syllable types occurred with varying frequencies (Table~\ref{tab:types}). Down FM was the most common ($\sim$25--30\%), followed by flat ($\sim$20--25\%) and short ($\sim$15--20\%). Complex and two-syllable types were rare ($<$5\% combined).

\begin{table}[h]
\centering
\caption{USV syllable type distribution.}
\label{tab:types}
\begin{tabular}{lc}
\toprule
Syllable Type & Approximate \% of Total \\
\midrule
Down FM & 25--30\% \\
Flat & 20--25\% \\
Short & 15--20\% \\
Chevron & 8--12\% \\
Step Down & 5--10\% \\
Multi & 5--8\% \\
Step Up & 3--5\% \\
Up FM & 3--5\% \\
Complex & 2--3\% \\
Two-syllable & 1--2\% \\
\bottomrule
\end{tabular}
\end{table}

\subsection{Two Intonation Mechanisms}

Analysis of the correlation between instantaneous airflow and pitch revealed two distinct intonation mechanisms (Table~\ref{tab:intonation}). Six of 10 syllable types showed predominantly positive intonation, where pitch changes tracked expiratory airflow changes ($r > 0$): down FM, flat, short, chevron, multi, and step down. In these syllable types, when airflow speed increased, pitch increased; when airflow speed decreased, pitch decreased.

Two syllable types showed predominantly negative intonation ($r < 0$): step up and up FM. In these types, pitch increased when airflow decreased, consistent with active laryngeal tension modulation independent of airflow. Complex and two-syllable types showed mixed intonation, using both mechanisms within single vocalizations.

\begin{table}[h]
\centering
\caption{Intonation classification by USV syllable type.}
\label{tab:intonation}
\begin{tabular}{llc}
\toprule
USV Type & Primary Intonation & Correlation Direction \\
\midrule
Down FM & Positive & $r > 0$ \\
Flat & Positive & $r > 0$ \\
Short & Positive & $r > 0$ \\
Chevron & Positive & $r > 0$ \\
Multi & Positive & $r > 0$ \\
Step Down & Positive/Mixed & $r \sim 0$ or mixed \\
Step Up & Negative & $r < 0$ \\
Up FM & Negative & $r < 0$ \\
Complex & Mixed & Both mechanisms \\
Two-syllable & Mixed & Both mechanisms \\
\bottomrule
\end{tabular}
\end{table}

\subsection{Diaphragm Re-Activation During Vocalization}

EMG recordings revealed the mechanism underlying positive intonation: re-activation of the diaphragm during the expiratory phase of vocalization breaths. During normal (non-USV) breathing, the diaphragm is active only during inspiration and silent during expiration. During USV breaths, the diaphragm showed an ectopic burst of activity during expiration, creating a brief inspiratory effort nested within the vocal breath.

This diaphragm re-activation slowed expiratory airflow (creating a dip in the airflow waveform), and the corresponding pitch decrease confirmed the positive intonation mechanism. When the diaphragm relaxed, expiratory airflow resumed and pitch increased. Laryngeal muscles (thyroarytenoid and cricothyroid) showed anti-phase activity with the diaphragm, suggesting coordinated but opposing modulation.

\subsection{iRO Optogenetic Activation Produces USVs}

Optogenetic activation of Penk+/Vglut2+ iRO neurons at 10 Hz entrained breathing to the stimulation frequency and increased respiratory amplitude. Critically, USVs were produced during the peak of entrained expirations. The evoked USVs included most of the 10 endogenous syllable types, with a strong bias toward positive intonation patterns. Negative intonation USVs were rare in the evoked repertoire, suggesting that the iRO primarily drives the breath-modulation mechanism, while negative intonation requires additional neural input from laryngeal premotor circuits.

Cell-type specificity was confirmed by comparing iRO activation with activation of neighboring neural populations: only Penk+/Vglut2+ neurons produced both robust breathing entrainment and USV production.

\subsection{Two-Mechanism Model of Intonation}

Based on these results, we propose a two-mechanism model of USV intonation (Table~\ref{tab:model}). Positive intonation is generated by breath waveform modulation via the iRO brainstem circuit, which controls diaphragm re-activation during expiration to shape expiratory airflow and thereby pitch. Negative intonation is generated by laryngeal tension modulation, presumably via the nucleus retroambiguus (RAm) and laryngeal premotor neurons. Complex pitch patterns arise from the combination of both mechanisms within single vocalizations.

\begin{table}[h]
\centering
\caption{Two-mechanism model of USV intonation.}
\label{tab:model}
\begin{tabular}{llll}
\toprule
Mechanism & Source & Neural Substrate & Effect \\
\midrule
Positive intonation & Breath modulation & iRO $\rightarrow$ breathing CPG & Pitch tracks airflow \\
Negative intonation & Laryngeal tension & RAm/laryngeal neurons & Pitch anticorrelates \\
Combined & Both & iRO + laryngeal & Complex patterns \\
\bottomrule
\end{tabular}
\end{table}

\section{Discussion}

Our findings challenge the prevailing view that laryngeal tension is the sole mechanism for pitch modulation in mammalian vocalizations. In mice, the majority of USV syllable types (6 of 10) use positive intonation, where pitch is controlled by the shape of the expiratory breath waveform rather than by laryngeal muscle dynamics. This breath-driven mechanism is generated by re-activation of the diaphragm during expiration, a motor pattern coordinated by the iRO brainstem CPG.

The existence of two distinct intonation mechanisms provides a new framework for understanding USV repertoire diversity. Rather than requiring independent laryngeal control for each syllable type, the mouse can generate most of its vocal repertoire through variations in a single breath-shaping mechanism, with laryngeal modulation adding the remaining types and contributing to complex combinations. This dual-mechanism architecture may explain why USV repertoires can be so diverse despite the relatively simple anatomy of the rodent larynx \citep{riede2011}.

The iRO's sufficiency for producing most USV types with positive intonation establishes it as a key node in the adult vocalization circuit. Previous work identified the iRO as a neonatal vocalization CPG \citep{wei2022}, and our results extend its role to adult USV production. The rarity of negative intonation in iRO-evoked USVs confirms that this mechanism depends on separate laryngeal premotor circuits.

The breath-driven mechanism has implications for understanding vocal communication more broadly. In birdsong, respiratory patterns play a critical role in shaping syllable structure \citep{goller2004,suthers2004}. Our results suggest that mammals may share a similar respiratory basis for vocal diversity, with breath shape contributing to pitch modulation in addition to the well-characterized laryngeal mechanisms. Whether breath-driven intonation exists in human speech---where respiratory patterns clearly affect prosody---is an intriguing question for future investigation.

The positive intonation mechanism also resolves the paradox noted in the introduction: while injecting air below the larynx increases pitch (consistent with the positive airflow-pitch correlation), natural USVs show variable airflow-pitch relationships because both positive and negative intonation mechanisms are active. When analyzed separately by syllable type, the airflow-pitch relationship becomes clear and consistent.

Several limitations should be noted. First, plethysmography measures whole-body pressure changes, which are a proxy for but not identical to subglottal pressure. Direct subglottal pressure measurements would provide more precise quantification of the airflow-pitch relationship. Second, the EMG recordings were chronic and may not have captured all relevant muscles. Third, the optogenetic experiments provide sufficiency but not necessity evidence for iRO involvement in natural USV production; loss-of-function experiments would complement these findings \citep{tschida2019}.

\section{Conclusion}

We have demonstrated that the shape of the expiratory breath waveform, controlled by diaphragm re-activation during expiration via the iRO brainstem circuit, is a primary mechanism for generating pitch patterns (intonation) in mouse ultrasonic vocalizations. This breath-driven mechanism accounts for the majority of the USV syllable repertoire, while laryngeal tension modulation generates the remaining types. The two-mechanism model provides a new framework for understanding vocal diversity in rodents and potentially in mammals more broadly.

\bibliographystyle{plainnat}
\bibliography{references}

\end{document}
